\subsection{Análisis}

La etapa de análisis dió comienzo con la definición de los \textbf{objetivos principales} del incremento:

\begin{itemize}
    \item \textbf{Lógica musical:} con la implementación de todos aquellos elementos de teoría musical necesarios para la creación de diagramas de acordes \textit{i.e. notas, acordes, instrumentos y diagramas de acordes}.
    \item \textbf{Algoritmo de cálculo de diagramas de acordes:} dado un instrumento y un \gls{acorde}, el algoritmo debe generar un diagrama del acorde que sea facilmente practicable y musicalmente correcto. 
\end{itemize}

Seguidamente, se seleccionaron los \textbf{requisitos funcionales} \tabref{tab:requisitosfuncionalesincremento2} relacionados con los objetivos anteriores.

\begin{table}[H]
    \centering
    \scriptsize
    \begin{tabular}{l p{0.58\textwidth} p{0.20\textwidth} l}
        \textbf{ID} & \textbf{Requisito Funcional} & \textbf{Funcionalidad} & \textbf{Prioridad} \\ \hline

        RF26 & El sistema debe aceptar distintos instrumentos de cuerda pulsada. & Cálculo de diagramas de acorde & Alta \\ \hline

        RF27 & El sistema debe aceptar como entrada uno o varios acordes. & Cálculo de diagramas de acorde & Alta \\ \hline

        RF28 & El sistema debe devolver un único diagrama de acorde por cada acorde solicitado. & Cálculo de diagramas de acorde & Alta \\ \hline

        RF29 & El sistema debe garantizar que cada diagrama generado sea musicalmente correcto respecto al acorde solicitado. & Cálculo de diagramas de acorde & Alta \\ \hline

        RF30 & El sistema debe garantizar que cada diagrama generado sea físicamente ejecutable. & Cálculo de diagramas de acorde & Alta \\ \hline

        RF31 & El sistema debe calcular cada diagrama optimizando simultáneamente la correctitud musical y la facilidad de ejecución. & Cálculo de diagramas de acorde & Alta \\ \hline

        RF32 & El sistema debe aceptar acordes en notación latina o anglosajona. & Cálculo de diagramas de acorde & Media \\ \hline

    \end{tabular}
    \caption{Requisitos funcionales del algoritmo de cálculo de diagramas de acordes}
    \label{tab:requisitosfuncionalesincremento2}
\end{table}

Finalmente, se realizó una investigación sobre los \textbf{fundamentos de teoría musical} \capref{c:fundamentos_musicales} y \textbf{instrumentos de cuerda pulsada} \capref{c:fundamentos_cuerda}, con el fín de conocer los detalles escenciales para la construcción y representación de diagramas de acordes.